\documentclass[a4paper,12pt]{article}
\usepackage{color}
\usepackage{graphicx, gensymb}
\usepackage{amsfonts, cancel, amssymb}
\usepackage{amsmath, amsthm, array, esvect, esint}
\usepackage{typearea}
\usepackage[a4paper,left=2cm,right=2cm,top=2cm,bottom=3cm,bindingoffset=5mm]{geometry}
\newcommand\numberthis{\addtocounter{equation}{1}\tag{\theequation}}
\newcommand{\bra}{}
\newcommand{\kom}{}
\newcommand{\brak}{}
\newcommand{\braket}{}
\newcommand{\intx}{}
\newcommand{\intr}{}
\newcommand{\kla}{}
\newcommand{\klam}{}
\newcommand{\klammer}{}
\newcommand{\oper}{}
\newcommand{\tecxt}{}
\newcommand{\Bra}{}
\newcommand{\Ket}{}

\begin{document}

\title{08 Dirac-Gleichung}
\author{Joachim Schwardt}
\date{\today}
\maketitle

\section*{Aufgabe 8.1: Dirac-Gleichung}
Seien $\sigma^j$ die Pauli-Matrizen $\sigma^1=\begin{pmatrix}
0&1\\ 1&0
\end{pmatrix}$, $\sigma^2=\begin{pmatrix}
0&-i\\ i&0
\end{pmatrix}$ und $\sigma^3=\begin{pmatrix}
1&0\\ 0&-1
\end{pmatrix}$. Dann definiert man die Dirac-Matrizen $\gamma^\mu$ über $\gamma^0=\begin{pmatrix}
I_2&0\\ 0&-I_2
\end{pmatrix}$ und $\gamma^j=\begin{pmatrix}
0&\sigma^j\\ -\sigma^j&0
\end{pmatrix}$. Betrachte die relativistische Energie-Impuls-Beziehung mit $p_\mu=\binom{E}{-\vec{p}}$:
\begin{align*}
m&=\sqrt{E^2-\vec{p}^2}=\sqrt{p^\mu p_\mu }\overset{!}{=}\alpha^\mu p_\mu&\Rightarrow &&p^\mu p_\mu\overset{!}{=}\alpha^\mu p_\mu\alpha^\nu p_\nu
\end{align*}
Dabei ist mit $g^{\mu\nu}=\begin{pmatrix}
1&0 \\ 0&-I_3
\end{pmatrix}$ die Minkowski-Metrik gemeint. Eine hinreichende Bedingung ist dabei $\klammer\left\{{\alpha^\mu,\alpha^\nu} \right\}=2g^{\mu\nu}$. Daraus folgt nämlich mit $\klam\left[ {p_\mu, p_\nu} \right]=0$:
\begin{align*}
2\alpha^\mu\alpha^\nu p_\mu p_\nu &=\alpha^\mu\alpha^\nu p_\mu p_\nu + \alpha^\nu\alpha^\mu p_\nu p_\mu=\klammer\left\{{\alpha^\mu,\alpha^\nu} \right\}p_\mu p_\nu=2g^{\mu\nu} p_\mu p_\nu=2m^2
\end{align*}
Die Dirac-Matrizen erfüllen gerade diese Anti-Kommutator-Relation. Die Dirac-Gleichung ergibt sich dann aus der Ersetzungsregel $p_\mu=i\partial_\mu=i\binom{\partial_t}{-\nabla}$:
\begin{align*}
m&=\gamma^\mu p_\mu &\Rightarrow && (i\gamma^\mu\partial_\mu - m)\psi&=0
\end{align*}
Betrachte die Wirkung der Dirac-Matrizen auf einen Spinor $u=(a,b,c,d)^\top$:
\begin{align*}
\gamma^\mu u&=\begin{pmatrix}
(a, b, -c, -d)^\top \\ (d, c, -b, -a)^\top \\
(-id, ic, ib, -ia)^\top \\ (c, -d, -a, b)^\top
\end{pmatrix}
\end{align*}
Die gegebenen Lösungen reproduzieren mit $0=\kla\left( {\gamma^\mu p_\mu -m} \right)u_1$ wieder $m^2=p^\mu p_\mu$. Dabei ist $u_1=(E+m,0,p_z,p_x+ip_y)^\top$:
\begin{align*}
mu_1&=N\begin{pmatrix}
E(E+m)&-p_x(p_x+ip_y)&+ip_y(p_x+ip_y)& -p_zp_z\\
 0&-p_xp_z& -ip_yp_z& +p_z(p_x+ip_y)\\
  -Ep_z&+0 &-0& +p_z(E+m) \\
   -E(p_x+ip_y)&+p_x(E+m)& +ip_y(E+m)&-0
\end{pmatrix}\\
\Rightarrow\ \ \ 0&=(E^2-\vec{p}^2,0,0,0)^\top
\end{align*}
\section*{Aufgabe 8.2: Nicht-relativistischer Grenzfall}
Es gilt $p_j\le |\vec{p}|=\gamma\beta m$ und $E=\gamma m$. Damit folgt für Terme der folgenden Form:
\begin{align*}
\frac{p_j}{E+m}&\le \frac{\gamma\beta m}{\gamma m +m}=\frac{\beta}{1+\sqrt{1-\beta^2}}\lesssim\beta
\end{align*}
\section*{Aufgabe 8.3: Spinor-Normierung}
Sei $\vec{p}=p_z\vec{e}_z$. Dann lauten die Lösungen:
\begin{align*}
(u_1,u_2,u_3,u_4)&=N\begin{pmatrix}
1&0&\frac{p_z}{E-m}&0\\
0&1&0&\frac{-p_z}{E-m}\\
\frac{p_z}{E+m}&0&1&0\\
0&\frac{-p_z}{E+m}&0&1
\end{pmatrix}
\end{align*}
Aus der Bedingung $u^\dagger u=2E$ kann man die Konstante $N$ mit $\frac{p^2}{E+m}=E-m$ bestimmen:
\begin{align*}
u^\dagger u&=N^2\kla\left( {1+\frac{p_z^2}{(E+m)^2}} \right)&\Rightarrow &&N&=\sqrt{\frac{2E}{1+\frac{E-m}{E+m}}}=\sqrt{E+m}
\end{align*}
Es gilt $S_zu_j=\frac{1}{2}(-1)^{j-1}u_j$, wobei $S_z=\frac{1}{2}\begin{pmatrix}
\sigma_z & 0\\ 0&\sigma_z
\end{pmatrix} $ ist. Dann gilt für den Hamiltonian $H=\vec{\alpha}\cdot\vec{p}+\beta m$ mit $\beta=\gamma^0$ und $\alpha^j=\gamma^0\gamma^j$ bei $p_x=p_y=0$:
\begin{align*}
S_zH&=\frac{1}{2}\begin{pmatrix}
\sigma_z & 0 \\ 0 & \sigma_z
\end{pmatrix}\begin{pmatrix}
m & p_z\sigma_z \\ 
p_z\sigma_z & -m
\end{pmatrix}=\frac{1}{2}\begin{pmatrix}
m\sigma_z & p_z \\ p_z & -m\sigma_z
\end{pmatrix}&\Rightarrow &&\klam\left[ {S_z,H} \right]&=0
\end{align*} 
\section*{Aufgabe 8.5: Chiralität und Helizität}
Es gilt $\gamma^5:=i\gamma^0\gamma^1\gamma^2\gamma^3\begin{pmatrix}
0&I_2 \\ I_2&0
\end{pmatrix} $. Damit sind die allgemeinen Lösungen keine Eigenvektoren von $\gamma^5$. Für $m=0$ ist aber $\frac{p_z}{E+m}=1$, womit es die Eigenwerte $\lambda=(+,-,+,-)$ für die $\vec{u}=(u_1,u_2,u_3,u_4)$ gibt ($+$ steht für rechtshändig). Für $p=p_z$ gilt $h:=\frac{\vec{S}\cdot\vec{p}}{|p|}=S_z$.
\section*{Aufgabe 8.6: Parität von Dirac-Fermionen}
Die Parität erhält man aus $\gamma^0u_\mu\overset{?}{=}c_\mu u_\mu$ zu $c_\mu = (+,+,-,-)^\top$. 
\end{document}